\documentclass[a4paper,12pt]{article}
\usepackage[utf8]{inputenc}
\usepackage{graphicx}
\usepackage{booktabs}
\usepackage{hyperref}
\usepackage{geometry}
\usepackage[english,ngerman]{babel}
\usepackage{amsmath}
\usepackage{amssymb}
\usepackage{float}
\usepackage{caption}

\usepackage[backend=biber]{biblatex}
%\usepackage{natbib}
\addbibresource{citations.bib}
\addbibresource{citations.bib}
\geometry{a4paper, margin=1in}

\begin{document}

\begin{titlepage}
    \centering
    {\Large Technische Universität Dortmund}\\[0.5cm]
    {\LARGE Fakultät Statistik}\\[1cm]

    \vspace*{5cm}
    {\Large {Erhebungstechniken}}\\[0.5cm]
    {\Large {Bericht}}\\[1.5cm]
    {\Huge \textbf{ (an ende machen) }}\\[1.5cm]
    {\Large Erheber von Dortmund}\\[3.5cm]
    { Dozenten }\\Prof. Dr. Philipp Doebler\\
    Dr. Susanne Frick\\
    Loreen Sabel, M.Sc.\\
    [0.5cm]
    { Gruppenmitglieder}\\Anastasiia Romanenko \\ Elaha Bahir\\ Furkan Koca\\[1.5cm]
    
    {\large 19.01.2026}
    \vfill
\end{titlepage}

\begingroup
  \flushbottom % 
  %\setlength{\parskip}{0pt plus 1fil}%
  \tableofcontents
  \newpage
\endgroup

\section{Einleitung}

Der Lernerfolg von Studierenden wird durch eine Vielzahl von Faktoren beeinflusst, wobei die bereitgestellten Lehr- und Lernmaterialien eine zentrale Rolle einnehmen. In der heutigen Hochschullehre stehen Studierende vor einer großen Auswahl an Materialien, darunter Skripte, Übungsaufgaben, Vorlesungsaufzeichnungen, Bücher sowie digitale und KI-gestützte Angebote. Trotz dieser Vielfalt bleibt häufig unklar, welche Materialien tatsächlich als lernförderlich wahrgenommen werden und wie sie den individuellen Lernerfolg beeinflussen.

Vor diesem Hintergrund verfolgt die vorliegende Studie das Ziel, zu untersuchen, wie Studierende der Technischen Universität Dortmund die ihnen zur Verfügung gestellten Lehr- und Lernmaterialien wahrnehmen und in welchem Zusammenhang diese mit dem subjektiv eingeschätzten Lernerfolg stehen. Dabei wird nicht nur betrachtet, welche Arten von Lernmaterialien genutzt werden, sondern auch, welche Eigenschaften und Rahmenbedingungen, wie Struktur, Verständlichkeit, Verfügbarkeit oder Verbindlichkeit, aus Sicht der Studierenden als besonders relevant empfunden werden.

Ein weiterer Fokus der Untersuchung liegt auf motivationalen und organisatorischen Aspekten des Lernens. So wird analysiert, inwiefern verpflichtende Übungsaufgaben, Fristen oder klar strukturierte Materialien das Lernverhalten beeinflussen und welche Rolle Lernmaterialien in Bezug auf Motivation, Prüfungsstress und selbstständiges Lernen spielen. Darüber hinaus werden individuelle Studienmerkmale, wie Fachsemester, Fakultätszugehörigkeit und angestrebter Abschluss, berücksichtigt, um mögliche Zusammenhänge differenzierter betrachten zu können.

Zur Beantwortung dieser Fragestellungen wurde eine empirische Studie in Form eines standardisierten Online-Fragebogens durchgeführt. Diese Methode ermöglicht es, die Einschätzungen und Erfahrungen einer größeren Gruppe von Studierenden systematisch zu erfassen und quantitativ auszuwerten. Der Fragebogen umfasst unter anderem Aspekte zur Nutzungshäufigkeit verschiedener Lernmaterialien, zur wahrgenommenen Qualität und Verfügbarkeit sowie zu deren Einfluss auf Lernerfolg, Motivation und Prüfungsbelastung.

Zunächst wird die Motivation der Studie dargelegt und die zugrunde liegenden Forschungsfragen formuliert. Anschließend wird das verwendete Erhebungsinstrument vorgestellt, wobei der Aufbau des Online-Fragebogens sowie die Gründe für dessen Einsatz erläutert werden. Darauf aufbauend werden die Stichprobe und der Datensatz beschrieben, einschließlich der Art der erhobenen Variablen und der vorgenommenen Datenaufbereitung. Im Ergebnisteil werden die zentralen Befunde der Untersuchung mithilfe von Grafiken und Tabellen dargestellt und inhaltlich erläutert, wobei sowohl deskriptive Auswertungen als auch explorative Analysen berücksichtigt werden. Abschließend werden die Ergebnisse im Diskussionsteil im Hinblick auf die Forschungsfragen interpretiert, Limitationen der Studie reflektiert und mögliche Verbesserungen der fragebogengestützten Erhebung diskutiert. 

\section{Problemstellung}

\section{Erhebungsinstrument}

\subsection{Stichprobe und Datensatz}

Der für diese Analyse verwendete Datensatz wurde von unseren Dozent:innen bereitgestellt. Die Daten stammen aus der elften Runde des \texttt{European Social Survey} (ESS11), einer international vergleichenden, akademisch geleiteten Umfrage. ESS11 umfasst Informationen aus 28 europäischen Staaten.

Die Datenerhebung erfolgte durch etwa einstündige, persönliche Interviews mit Personen ab 15 Jahren, die in privaten Haushalten leben – unabhängig von Nationalität, Staatsbürgerschaft, Sprache oder rechtlichem Status. Die Erhebungsphase der elften Runde erstreckte sich vom 8. März 2023 bis zum 4. Juli 2024.

Der Datensatz umfasst 46\,162 Beobachtungen und siebzehn Variablen. Diese beinhalten sowohl personenbezogene Merkmale – wie Alter, Geschlecht und eine etwaige COVID-19-Erkrankung – als auch Angaben zum Vertrauen der Befragten in politische Institutionen. Einen detaillierten Überblick über die numerischen Merkmale bietet Tabelle~\ref{tab:features_overview_blatt}. Zur besseren Verständlichkeit wurden die ursprünglichen Variablennamen umbenannt; in der Analyse wird durchgehend die in der Tabelle angegebene Bezeichnung verwendet.

\begin{minipage}{\textwidth}
    \centering
    \begin{tabular}{l p{8cm} p{4cm}}
        \toprule
        Merkmal & Beschreibung & Wertebereich \\
        \midrule
        \texttt{person\_id} & Eindeutige ID-Nummer der befragten Person & [50002, 211884] \\
        \texttt{covid\_infection} & Selbstberichtete COVID-19-Infektion (1 = positiver Test, 2 = vermutete Infektion, 3 = keine Infektion) & \{1, 2, 3\} \\
        \texttt{vaccined} & COVID-19-Impfstatus: mindestens eine Dosis eines in [Land] zugelassenen Impfstoffs erhalten (1 = ja, 2 = nein) & \{1, 2\} \\ 
        \texttt{education\_niveau} & Bildungsniveau & \{1, \dots, 7, 55\} \\
        \texttt{house\_income} & Haushaltseinkommen (Dezile) & \{1, \dots, 10\} \\
        \texttt{houshold\_n} & Anzahl der Personen im Haushalt & \{0, \dots, 13\} \\
        \texttt{internet\_frequency} & Häufigkeit der Internetnutzung (1 = nie bis 5 = täglich) & \{1, \dots, 5\} \\
        \texttt{gender} & Geschlecht (1 = männlich, 2 = weiblich) & \{1, 2\} \\
        \texttt{marital\_status} & Familienstand & \{1, \dots, 6\} \\
        \texttt{residency\_type} & Art des Wohnorts (von Großstadt bis ländlich) & \{1, \dots, 5\} \\
        \texttt{age} & Alter der befragten Person & [15, 90] \\
        \texttt{trust\_parlament} & Vertrauen in das nationale Parlament & \{1, \dots, 10\} \\
        \texttt{trust\_legalsys} & Vertrauen in das Rechtssystem & \{1, \dots, 10\} \\
        \texttt{trust\_police} & Vertrauen in die Polizei & \{1, \dots, 10\} \\
        \texttt{trust\_politics} & Vertrauen in Politiker:innen & \{1, \dots, 10\} \\
        \texttt{trust\_partys} & Vertrauen in politische Parteien & \{1, \dots, 10\} \\
        \bottomrule
    \end{tabular}
    \captionof{table}{Detaillierter Überblick über die im Datensatz enthaltenen Variablen}
    \label{tab:features_overview_blatt}
\end{minipage}

\vspace{1em}



Einige Variablen erfordern eine genauere Erläuterung hinsichtlich ihrer Kodierung. Die Variable \texttt{education\_niveau} basiert auf dem internationalen Standard ISCED (International Standard Classification of Education) und nimmt Werte von eins bis sieben an. Diese können als aufsteigende Stufen formaler Bildungsabschlüsse interpretiert werden – je höher der Wert, desto höher der Bildungsgrad.

Auch die Variable \texttt{marital\_status}, die den Familienstand abbildet, ist in sechs Kategorien unterteilt: (1) gesetzlich verheiratet, (2) in einer gesetzlich eingetragenen Partnerschaft, (3) gesetzlich getrennt lebend, (4) gesetzlich geschieden, (5) verwitwet und (6) keine der genannten Optionen.


Die sogenannten Vertrauensvariablen erfassen das Vertrauen der Befragten in verschiedene staatliche Institutionen. Hierzu wurden die Teilnehmenden gebeten, ihr Vertrauen in jede Institution auf einer Skala von 1 (sehr geringes Vertrauen) bis 10 (sehr hohes Vertrauen) einzuschätzen.

Zudem enthält der Datensatz die Variable \texttt{country}, die angibt, in welchem Land die jeweilige Person befragt wurde. Diese Variable umfasst 28 Ausprägungen, wobei jeweils der standardisierte Länderkürzel verwendet wird.


Zusätzlich wurde in der Umfrage ein Kodierungssystem für Antwortverweigerungen oder fehlende Angaben verwendet. Diese speziellen Werte kennzeichnen Fälle, in denen eine Person eine Frage nicht beantworten wollte oder konnte. Da solche Fälle im Datensatz nur sehr selten auftreten, wurden diese Werte in der Spalte Abbildermenge der Übersichtstabelle nicht aufgeführt. Im Kapitel~4.1 wird der Umgang mit den fehlenden Werten näher erläutert.


\subsection{Datenqualität}

Die Qualität der im ESS11 erhobenen Daten ist grundsätzlich hoch, dennoch müssen einige Aspekte berücksichtigt werden, die die Vergleichbarkeit und Aussagekraft beeinflussen können.

Zunächst unterscheiden sich die Stichprobendesigns zwischen den teilnehmenden Län- \\ dern teils erheblich. Ziel dieser Designs ist es, in jedem Land eine vergleichbare effektive Stichprobengröße zu erreichen, um möglichst ähnliche Konfidenzintervalle gewährleisten zu können. In  vielen Ländern basiert die Auswahl auf einer Adressstichprobe: Dabei wird an einer zufällig ausgewählten Adresse eine Person zufällig ausgewählt. Diese Methode führt allerdings dazu, dass Personen, die allein in einem Haushalt leben, eine höhere Auswahlwahrscheinlichkeit haben als Personen in Mehrpersonenhaushalten. Dies kann die Repräsentativität verzerren, insbesondere in Bezug auf Merkmale, die mit Haushaltsgröße oder -struktur korrelieren.

Einflussfaktoren auf die Datenqualität ergeben sich unter anderem aus Unterschieden im Umgang mit Nonresponse zwischen den Ländern, etwa in der Nachbearbeitung, Schulung des Erhebungspersonals oder Teilnahmebereitschaft. Solche Unterschiede können zu systematischen Verzerrungen führen. Zudem weisen einige Variablen, wie \texttt{house\_income}, viele fehlende Werte auf. Auch die Zahl der Beobachtungen variiert je nach Land, was bei der Interpretation der Ergebnisse berücksichtigt werden sollte.


\section{Ergebnisse}

Dieses Kapitel liefert die Anwendung der beschriebenen Methoden und die Auswertung der Ergebnisse. Die Daten wurden bearbeitet mithilfe der Software Python \cite{python} und Python-Bibliotheken pandas \cite{reback2020pandas}, plotly \cite{plotly2015}, scipy \cite{virtanen2020scipy}, matplotlib \cite{hunter2007matplotlib}, statsmodels \cite{seabold2010statsmodels} und seaborn \cite{seabold2010statsmodels}.

\subsection{quasi demografie}


\subsection{Ergebnisse zur Nutzung und Bewertung von Lernmaterialien}

In diesem Unterkapitel geben wir einen deskriptiven Überblick über die Antworten auf die wichtigsten Fragen des Fragebogens und untersuchen, welche Lernmaterialien den Studierenden wichtig sind und wie sie genutzt werden. Zu diesem Zweck wurden gestapelte Balkendiagramme für alle kategorialen Variablen erstellt. Die entsprechenden Visualisierungen sind in Abbildung~\ref{fig:materials_nutzung} und Abbildung~\ref{fig:materials_qualities} dargestellt.

Zur Verbesserung der Verständlichkeit und Lesbarkeit werden die Variablen in diesem Abschnitt nach ihrem Durchschnitt sortiert und von oben nach unten dargestellt. Jeder Balken steht für ein Lernmaterial, die einzelnen Segmente zeigen die prozentualen Anteile der fünf Antwortkategorien. Die Prozentwerte ergeben zusammen jeweils 100\,\% pro Lernmaterial.

Abbildung~\ref{fig:materials_nutzung} zeigt die Nutzung verschiedener Lernmaterialien durch Studierende. Aus der Abbildung wird deutlich, dass sich die Nutzung der Lernmaterialien teilweise stark unterscheidet. Besonders häufig werden Übungsaufgaben mit Musterlösung sowie Altklausuren verwendet. Bei diesen Materialien entfällt ein großer Anteil auf die höchsten Nutzungskategorien, was darauf hindeutet, dass sie für die Prüfungsvorbereitung eine wichtige Rolle spielen. Auch KI-Chatbots werden von vielen Studierenden regelmäßig genutzt.

\begin{figure}[htpb]
    \centering
    \includegraphics[width=16cm]{plots/learning_materials_group_sorted.pdf}
    \caption{Gestapeltes Balkendiagramm der von Studierenden genutzten Lernmaterialien}
    \label{fig:materials_nutzung}
\end{figure}

Lernmaterialien wie Vorlesungsvideos, Online-Videos, Skripte und Bücher werden überwiegend im mittleren Nutzungsbereich eingeordnet. Dies spricht dafür, dass sie zwar regelmäßig verwendet werden, jedoch nicht für alle Studierenden die wichtigste Lernquelle darstellen. Bücher weisen dabei den größten Anteil an niedrigen Bewertungen (Kategorie 1) auf, liegen im Mittelwert jedoch nicht auf dem letzten Platz.

Insgesamt zeigt die Abbildung, dass Studierende vor allem solche Lernmaterialien bevorzugen, die einen direkten Bezug zu Übungen und Prüfungen haben, während Materialien, die eher der Erweiterung des Wissens dienen, eine untergeordnete Rolle spielen.

Abbildung~\ref{fig:materials_qualities} zeigt die Einschätzung der Studierenden zu verschiedenen Aussagen über die von ihnen genutzten Lernmaterialien.

\begin{figure}[htpb]
    \centering
    \includegraphics[width=16cm]{plots/materials_that_i_use_sorted.pdf}
    \caption{Gestapeltes Balkendiagramm wichtiger Eigenschaften von Lernmaterialien aus Sicht der Studierenden}
    \label{fig:materials_qualities}
\end{figure}

Die Ergebnisse zeigen insgesamt eine eher positive Bewertung der Lernmaterialien. Viele Aussagen weisen hohe Anteile in den oberen Skalenpunkten (4 und 5) auf. Besonders deutlich wird dies bei Aussagen wie „geben mir Sicherheit, den Stoff zu schaffen“, „fördern Verstehen statt Auswendiglernen“ sowie „sind klar strukturiert und organisiert“. Auch die Unterstützung des selbstständigen Lernens und die Erleichterung des Lernprozesses werden überwiegend positiv bewertet.

Ein gemischteres Bild zeigt sich bei den Aussagen zum Prüfungsstress und zur Motivation. Zwar geben einige Studierende an, dass die Lernmaterialien ihren Prüfungsstress reduzieren und die Motivation steigern, jedoch liegen hier auch viele Antworten im mittleren Skalenbereich.

Hinsichtlich des Zeitaufwands und der Arbeitsbelastung zeigen die Ergebnisse, dass ein Teil der Studierenden angibt, durch die Lernmaterialien einen erhöhten Zeitaufwand und eine höhere Arbeitsbelastung zu erleben. Gleichzeitig werden den Materialien jedoch häufig eine gute Struktur, hohe Verständlichkeit sowie eine klare Prüfungsrelevanz zugeschrieben.


\subsection{nebenfragestellung}




\section{Diskussion}

In der vorliegenden Analyse wurde untersucht, welche Faktoren die Entscheidung zur COVID-19-Impfung beeinflussen. Dazu wurden Daten der elften Runde des European Social Survey (ESS11) herangezogen, die umfangreiche Informationen aus 28 europäischen Ländern enthalten. Die Analyse basierte auf einem Datensatz mit über 46.000 Beobachtungen, wobei personenbezogene Merkmale, sozioökonomische Indikatoren sowie Vertrauen in staatliche Institutionen als Prädiktoren berücksichtigt wurden.

Zur Vorbereitung der Daten wurde eine umfangreiche Datenbereinigung und - transformation durchgeführt. Darüber hinaus wurden die Annahmen der logistischen Regression geprüft und bei einigen Variablen durch Quadratwurzel-Transformationen korrigiert. Zur Berücksichtigung der Klassenungleichgewichte wurde eine Gewichtung der Zielvariable implementiert.

Die Analyse ergab, dass vor allem Alter, Bildungsniveau und Haushaltseinkommen starke positive Prädiktoren für die Impfwahrscheinlichkeit sind. Ältere Personen sowie solche mit höherer Bildung und höherem Einkommen zeigen deutlich erhöhte Impfquoten. Ebenso spielte das Vertrauen in staatliche Institutionen, insbesondere in Parlament, Rechtssystem und Polizei, eine wichtige Rolle und erhöhte die Impfbereitschaft signifikant. Darüber hinaus zeigten sich ausgeprägte länderspezifische Unterschiede in der Impfquote, wobei einige südeuropäische Länder hohe Impfraten aufwiesen, während einige osteuropäische Länder vergleichsweise niedrige Werte zeigten. Weitere Faktoren wie Geschlecht, Familienstand oder Erfahrungen mit einer COVID-19-Infektion beeinflussten die Entscheidung ebenfalls signifikant, allerdings mit teils moderaterer Stärke. Interessanterweise hatte die tägliche Internetnutzung keinen signifikanten Einfluss.

Trotz der umfangreichen Daten und der sorgfältigen Analyse weist die Studie einige Einschränkungen auf. Die Stichprobenauswahl und das Erhebungsdesign unterscheiden sich zwischen den Ländern, was die Vergleichbarkeit einschränken und potenzielle Verzerrungen erzeugen kann, etwa durch unterschiedliche Nonresponse-Raten oder Haushaltsgrößen. Zudem wurden fehlende Werte insbesondere beim Haushaltseinkommen durch eine Indikatorvariable ersetzt, was zwar den Verlust von Beobachtungen minimiert, aber möglicherweise Verzerrungen in den Schätzungen hervorruft. Die Vereinfachung kategorialer Variablen zu kontinuierlichen Variablen kann ebenfalls zu Informationsverlust führen. Außerdem kann die logistische Regression nur lineare Effekte erfassen; komplexere Zusammenhänge oder Interaktionen bleiben unberücksichtigt. Nicht zuletzt besteht die Möglichkeit, dass wichtige Einflussfaktoren außerhalb des ESS-Datensatzes liegen, die hier nicht erfasst wurden.

Insgesamt liefert die Analyse wertvolle Einblicke in die Determinanten der COVID-19-Impfbereitschaft in Europa, betont aber zugleich die Notwendigkeit weiterer Forschung mit noch differenzierteren Daten und Methoden, um die komplexen Zusammenhänge umfassender abzubilden.

\section{Reflexion der fragebogengestützten Erhebung}

\newpage

%\section{Literaturverzeichnis}
%\nocite{*}
\nocite{*}

\newpage

\appendix

\end{document}
